\subsection{Operator Polynomials and Evaluation}

% **Role:** structure revealed by invariant subspaces and eigenbehavior.
%
% Characteristic polynomial, minimal polynomial, Cayley--Hamilton, generalized eigenvectors, nilpotent parts.

Polynomial structure records the equations an operator satisfies and turns them into invariant pieces through characteristic and minimal polynomials, Cayley--Hamilton, generalized eigenspaces, and nilpotent parts.

\begin{definition} \label{def:operator-powers} \label{def:polynomial-applied-to-operator} \label{def:annihilating-polynomial} \hlt{Operator Powers and Polynomial Evaluation}\\
Let $T\in\LL(\mathcal{V})$.
\begin{enumerate}[label=(\roman*)]
\setlength{\itemsep}{0pt}
\item The \hlt{powers} of $T$ are $T^m\defeq\underbrace{T\cdots T}_{m\text{ times}}$ for $m\in\mathbb{N}$, with $T^0\defeq I_{\mathcal{V}}$. If $T$ is invertible, set $T^{-m}\defeq(T^{-1})^m$.
\item For $p(z)=\sum\limits_{j=0}^{m}a_jz^j\in\mathbb{F}[z]$, the \hlt{polynomial $p$ applied to $T$} is $p(T)\defeq\sum\limits_{j=0}^{m}a_jT^j$.
\item $p\in\mathbb{F}[z]$ is an \hlt{annihilating polynomial} for $T$ if $p(T)=0_{\LL(\mathcal{V})}$; equivalently, $p$ \hlt{annihilates} $T$.
\end{enumerate}
\end{definition}

\begin{definition} \label{def:product-of-polynomials} \hlt{Product of Polynomials}\\
If $p,q\in\mathbb{F}[z]$, their product $pq\in\mathbb{F}[z]$ is defined by $(pq)(z)=p(z)q(z)$ for $z\in\mathbb{F}$.
\end{definition}

\begin{proposition} \label{prop:operator-polynomial-products} \hlt{Multiplicative Properties of Operator Polynomials}\\
Let $p,q\in\mathbb{F}[z]$ and $T\in\LL(\mathcal{V})$. Then
\begin{enumerate}[label=(\roman*)]
\setlength{\itemsep}{0pt}
\item $(pq)(T)=p(T)q(T)$;
\item $p(T)q(T)=q(T)p(T)$.
\end{enumerate}
\begin{proof}
\begin{enumerate}[label=(\roman*)]
\setlength{\itemsep}{0pt}
\item Write $p(z)=\sum_{j=0}^{m}a_jz^j$ and $q(z)=\sum_{k=0}^{n}b_kz^k$. Then $(pq)(z)=\sum\limits_{j=0}^{m}\sum\limits_{k=0}^{n}a_jb_kz^{j+k}$.\\
Substituting $T$ gives $(pq)(T)=\sum\limits_{j=0}^{m}\sum\limits_{k=0}^{n}a_jb_kT^{j+k} =\left(\sum\limits_{j=0}^{m}a_jT^j\right)\left(\sum\limits_{k=0}^{n}b_kT^k\right) =p(T)q(T)$.
\item Since polynomial multiplication is commutative, $pq=qp$.\\
By part (i), $p(T)q(T)=(pq)(T)=(qp)(T)=q(T)p(T)$.
\end{enumerate}
\hlt{Structural Commentary}: Polynomial evaluation preserves products.\\
\hlt{Key Insight}: Expanding a polynomial product works the same with $T$ as with $z$.
\end{proof}
\end{proposition}

\begin{proposition} \label{prop:finite-dimensional-operators-have-annihilating-polynomials} \hlt{Existence of Annihilating Polynomials}\\
Let $\mathcal{V}$ be finite-dimensional with $n=\dim(\mathcal{V})>0$, and let $T\in\LL(\mathcal{V})$.\\
Then $T$ has a nonconstant annihilating polynomial of degree at most $n^2$.
\begin{proof}
By Theorem \ref{thm:linear-maps-and-matrices-same-data}, the vector space $\LL(\mathcal{V})$ has dimension $n^2$, so the $n^2+1$ operators $I_{\mathcal{V}},T,T^2,\ldots,T^{n^2}$ are linearly dependent. Hence there are scalars $a_0,\ldots,a_{n^2}$, not all zero, such that $\sum_{j=0}^{n^2}a_jT^j=0$.\\
If $a_1=\cdots=a_{n^2}=0$, then $a_0I_{\mathcal{V}}=0$, forcing $a_0=0$ because $\mathcal{V}\neq\{\mathbf{0}\}$, a contradiction. Thus $p(z)=\sum_{j=0}^{n^2}a_jz^j$ is a nonconstant annihilating polynomial.\\
\hlt{Structural Commentary}: Finite-dimensionality makes the powers of one operator linearly dependent inside $\LL(\mathcal{V})$.\\
\hlt{Key Insight}: A dependence among $I_{\mathcal{V}},T,\ldots,T^{n^2}$ is exactly a polynomial equation satisfied by $T$.
\end{proof}
\end{proposition}

\begin{proposition} \label{prop:polynomial-eigenvector-evaluation} \hlt{Polynomial Evaluation on Eigenvectors}\\
Let $T\in\LL(\mathcal{V})$, let $p\in\mathbb{F}[z]$, and suppose $T\mathbf{v}=\lambda\mathbf{v}$ with $\mathbf{v}\neq\mathbf{0}$. Then $p(T)\mathbf{v}=p(\lambda)\mathbf{v}$.
\begin{proof}
Induct on $j$ to get $T^j\mathbf{v}=\lambda^j\mathbf{v}$ for every $j\geq0$.\\
If $p(z)=\sum\limits_{j=0}^{m}a_jz^j$, then $p(T)\mathbf{v} =\sum\limits_{j=0}^{m}a_jT^j\mathbf{v} =\sum\limits_{j=0}^{m}a_j\lambda^j\mathbf{v} =p(\lambda)\mathbf{v}$.\\
\hlt{Structural Commentary}: On a one-dimensional eigenspace, $T$ acts like the scalar $\lambda$.\\
\hlt{Key Insight}: Polynomial expressions in $T$ reduce to the same polynomial evaluated at $\lambda$.
\end{proof}
\end{proposition}

